\chapter{Authenticity in the random oracle model}
\label{chap:random-oracle}

We work in the random oracle model \cite{BR93}. The hash is now a uniformly random function $H \colon \{0, \dots, n - 1\} \times V \to V$, where
$n$ is the number of nodes, and $\Delta$ is uniform on $V \setminus \{0\}$. The adversary sees
the material, makes $q$ adaptive queries to $H$, and outputs a label. The active labels of the
inputs and of the constants are first fixed parameters; Theorem~\ref{thm:garble-rom} at the end of
the chapter removes this restriction and bounds the real garbling with uniform zero-labels.

\section{The trace of the garbler}

\begin{definition}[Trace]
  \label{def:trace}
  \lean{GcZk.RO.idealWiresAux, GcZk.RO.nodeMux, GcZk.RO.muxAt}
  \leanok
  \uses{def:ideal}
  The trace records, for each position $i$ that holds a multiplexer, its active selector label
  $S_i$ and the XOR $c_i$ of its candidate values. The garbler reads $H$ at $(i, S_i)$ and the pad
  at $(i, S_i, c_i)$. In the real scheme the \emph{inactive point} of position $i$ is
  $(i, S_i + \Delta)$.
\end{definition}

\begin{lemma}[Congruence]
  \label{lem:walk-congr}
  \lean{GcZk.RO.muxAt_lt, GcZk.RO.evalAux_ideal, GcZk.RO.node_congr, GcZk.RO.walk_congr}
  \leanok
  \uses{def:trace}
  Two runs of the ideal garbler whose hashes and pads agree at the points named by the trace of
  the first produce the same material, the same wires and the same trace.
\end{lemma}

\begin{proof}
  \leanok
  Induction on the nodes; a node reads only the hash and the pad of its own position.
\end{proof}

\section{Adversaries and the two games}

\begin{definition}[Adversary]
  \label{def:qadv}
  \lean{GcZk.RO.QAdv, GcZk.RO.answers}
  \leanok
  A deterministic adversary chooses each query from the material and the previous answers, and its
  output from the material and all $q$ answers. A randomized adversary is a family of these
  indexed by its coins.
\end{definition}

\begin{lemma}[Identical until bad]
  \label{lem:answers-congr}
  \lean{GcZk.RO.answers_congr}
  \leanok
  \uses{def:qadv}
  Two oracles that agree at every query the adversary asks give the same answers.
\end{lemma}

This is the deterministic core of the fundamental lemma of game playing of Bellare and Rogaway
\cite[Lemma~2]{BR06}; see also Shoup \cite{Sho04}. In the proofs below it is used through injections
between sets of randomness rather than through probabilities.

\begin{proof}
  \leanok
  Induction on the number of queries.
\end{proof}

\begin{definition}[Real and simulated games]
  \label{def:games}
  \lean{GcZk.RO.extH, GcZk.RO.extK, GcZk.RO.rpad, GcZk.RO.realMat, GcZk.RO.RealWin,
    GcZk.RO.simMat, GcZk.RO.simQuery, GcZk.RO.SimWin, GcZk.RO.SimBad}
  \leanok
  \uses{def:trace, def:qadv, lem:auth-win-iff}
  In the real game the material uses the pads $H(i, S + \Delta) + c\,\Delta$ and the adversary
  wins if it outputs $A + \Delta$. In the simulated game the material uses independent pads
  $K \in V^n$, so nothing the adversary sees depends on $\Delta$; it wins with the same condition,
  and it hits if one of its queries is an inactive point $(i, S_i + \Delta)$ of the simulated
  trace.
\end{definition}

\section{The reduction}

\begin{lemma}[Moving the inactive values into the pads]
  \label{lem:to-sim}
  \lean{GcZk.RO.realMux, GcZk.RO.moveInactive, GcZk.RO.toSim, GcZk.RO.toSim_walk,
    GcZk.RO.moveInactive_injective, GcZk.RO.toSim_injective}
  \leanok
  \uses{lem:walk-congr, def:games}
  Fix $\Delta \neq 0$. Given $(H, J)$ with $J \in V^n$, let $H'$ be $H$ with its value at each
  inactive point $(i, S_i + \Delta)$ replaced by $J_i$, and let $K_i = H(i, S_i + \Delta) + c_i\Delta$
  at multiplexers and $K_i = J_i$ elsewhere. Then the simulated game with $(H', K)$ has the same
  material, wires and trace as the real game with $H$, and the map $(H, J) \mapsto (H', K)$ is
  injective.
\end{lemma}

\begin{proof}
  \leanok
  Active and inactive points differ because $\Delta \neq 0$, so $H'$ agrees with $H$ at the active
  points, and the pads agree by construction; Lemma~\ref{lem:walk-congr} gives the first claim.
  Both traces coincide, so $H$ is recovered from $H'$ off the inactive points and from $K$ at
  them, and $J$ from $H'$ at the inactive points and from $K$ elsewhere.
\end{proof}

\begin{lemma}
  \label{lem:real-win-to-sim}
  \lean{GcZk.RO.activeOut_ideal, GcZk.RO.realWin_toSim}
  \leanok
  \uses{lem:to-sim, lem:answers-congr}
  If the adversary wins the real game with $H$, it wins or hits in the simulated game with
  $(H', K)$.
\end{lemma}

\begin{proof}
  \leanok
  If it does not hit, its queries avoid the inactive points, where alone $H$ and $H'$ differ;
  by Lemma~\ref{lem:answers-congr} it gets the same answers, hence outputs the same label.
\end{proof}

\begin{lemma}
  \label{lem:card-sim}
  \lean{GcZk.RO.card_simWin_or_bad}
  \leanok
  \uses{def:games}
  For fixed $(H', K)$, at most $q + 1$ values of $\Delta$ give a simulated win or a hit.
\end{lemma}

\begin{proof}
  \leanok
  A win fixes $\Delta = Z - A$, and the $k$-th query $(i, y)$ hits only for $\Delta = y - S_i$.
\end{proof}

\begin{theorem}[Authenticity in the random oracle model]
  \label{thm:rom-bound}
  \lean{GcZk.RO.rom_bound}
  \leanok
  \uses{lem:real-win-to-sim, lem:card-sim}
  Let $V$ be finite. Summed over $\Delta \neq 0$, the hash functions for which a $q$-query
  adversary outputs the inactive output label number at most $(q + 1)\,|V|^{n |V|}$. With
  $\Delta$ uniform on $V \setminus \{0\}$ and $H$ uniform, the success probability is at most
  \[ \frac{q + 1}{|V| - 1} . \]
\end{theorem}

\begin{proof}
  \leanok
  For each $\Delta$, Lemmas~\ref{lem:to-sim} and \ref{lem:real-win-to-sim} inject the real wins,
  times $|V|^n$ choices of $J$, into the simulated wins or hits. Summing over $\Delta$ and swapping
  the sums, Lemma~\ref{lem:card-sim} bounds the total by $(q + 1)$ times the number of pairs
  $(H', K)$. Divide by $|V|^n$.
\end{proof}

\section{The real garbling}

\begin{lemma}
  \label{lem:const-congr}
  \lean{GcZk.Auth.idealAux_congr_const}
  \leanok
  \uses{def:ideal}
  The ideal garbler reads the active label of a constant only at the bit of that constant.
\end{lemma}

\begin{proof}
  \leanok
  Induction on the nodes.
\end{proof}

\begin{lemma}[Reparametrization]
  \label{lem:reparam}
  \lean{GcZk.Sec.extF, GcZk.Sec.seedOf, GcZk.Sec.activeIn, GcZk.Sec.shiftC, GcZk.Sec.constLab,
    GcZk.Sec.inpOf, GcZk.Sec.acOf, GcZk.Sec.gMat, GcZk.Sec.GarbleWin,
    GcZk.Sec.activeInputs_seedOf, GcZk.Sec.gMat_eq, GcZk.Sec.garbleWin_iff}
  \leanok
  \uses{thm:real-ideal, lem:const-congr, def:games}
  Let the seed consist of $\Delta$, zero-labels $L^{\mathrm{in}}$ of the inputs $j < k$ and
  zero-labels $L^{\mathrm{c}}$ of the constants, and let $x_j = 0$ for $j \geq k$. Then the
  material of the real garbling, and the game in which the adversary also receives the active
  labels of $x$, coincide with the game of Theorem~\ref{thm:rom-bound} with the fixed active labels
  $L^{\mathrm{in}}_j + x_j\Delta$ and $L^{\mathrm{c}}_i + b_i \Delta$.
\end{lemma}

\begin{proof}
  \leanok
  Theorem~\ref{thm:real-ideal} and Lemma~\ref{lem:const-congr}.
\end{proof}

\begin{theorem}[Authenticity of the real garbling]
  \label{thm:garble-rom}
  \lean{GcZk.Sec.garble_rom_bound}
  \leanok
  \uses{lem:reparam, thm:rom-bound}
  Let $V$ be finite and let the garbler choose $\Delta \neq 0$, the hash $H$ and the zero-labels
  of the inputs and of the constants uniformly. Let $x$ be an input fixed in advance, with
  $x_j = 0$ for $j \geq k$ as in Lemma~\ref{lem:reparam}. An adversary that sees the material and
  the active labels of $x$, and makes $q$ adaptive queries to $H$, outputs the inactive output label
  with probability at most $(q + 1)/(|V| - 1)$.
\end{theorem}

\begin{proof}
  \leanok
  For fixed $\Delta$, the maps $L^{\mathrm{in}} \mapsto L^{\mathrm{in}} + x\Delta$ and
  $L^{\mathrm{c}} \mapsto L^{\mathrm{c}} + b\Delta$ are bijections. By Lemma~\ref{lem:reparam}
  the sum over $\Delta$ and the zero-labels becomes a sum over the active labels, outside, of the
  sums of Theorem~\ref{thm:rom-bound}, inside.
\end{proof}

\begin{remark}
  The bound of Theorem~\ref{thm:rom-bound} does not use characteristic $2$: it holds for labels in
  any finite abelian group. Theorem~\ref{thm:garble-rom} does use it, through
  Theorem~\ref{thm:real-ideal}. With $128$-bit labels the authenticity advantage of an adversary
  making $q$ hash queries is at most $(q + 1)/(2^{128} - 1)$. Theorem~\ref{thm:sound} suggests the
  same bound for the soundness error of the protocol in the hybrid model; the formal statement,
  with the hash commitment and a prover with a bounded number of queries, is
  Theorem~\ref{thm:e2e}.
\end{remark}
